\usepackage{amsmath,amsthm}
\usepackage{tikz}

\usepackage{fontspec,unicode-math}
\setmainfont{texgyrepagella}[
  Extension=.otf,
  UprightFont=*-regular,
  BoldFont=*-bold,
  ItalicFont=*-italic,
  BoldItalicFont=*-bolditalic
]
\setmathfont{texgyrepagella-math.otf}

\numberwithin{equation}{section}

\usepackage{showkeys}

\newcommand{\MM}{\overline{\mathcal{M}}}
\newcommand{\EE}{\mathbb{E}}
\newcommand{\dd}{\mathrm{d}}

\DeclareMathOperator{\Cin}{\mathrm{Cin}}

% Keep the existing mathematical headings; optional exposition has roman bodies.
\newtheoremstyle{paperplain}
  {\topsep}{\topsep}{\itshape}{}{\bfseries}{}{.5em}
  {\thmname{#1}\thmnumber{ #2}\thmnote{ (#3)}}
\newtheoremstyle{paperremark}
  {\topsep}{\topsep}{\normalfont}{}{\bfseries}{}{.5em}
  {\thmname{#1}\thmnumber{ #2}\thmnote{ (#3)}}

\theoremstyle{paperplain}
\newtheorem{thm}{Theorem}[section]
\newtheorem{dfn}[thm]{Definition}
\newtheorem{lmm}[thm]{Lemma}
\newtheorem{prp}[thm]{Proposition}
\newtheorem{crl}[thm]{Corollary}

\theoremstyle{paperremark}
\newtheorem{rmk}{Remark}[section]
\newtheorem{exm}{Example}[section]
\theoremstyle{paperplain}

\newenvironment{prf}{\noindent \textbf{Proof: }}{\hfill $\mdlgwhtsquare$}
\newenvironment{prfn}[1]{\vskip 1em \noindent \textbf{Proof of #1: }}{\hfill $\mdlgwhtsquare$}
